\begin{resumo}
    A dengue é um dos principais desafios de saúde pública no Brasil, e a epidemia de 2024, com cerca de 6,5 milhões de notificações, evidenciou a necessidade de ferramentas de antecipação de surtos. Este trabalho desenvolveu e avaliou um modelo de previsão semanal de notificações de dengue para as 137 mesorregiões brasileiras, com quatro semanas de antecedência, a partir da integração dos microdados do SINAN com os dados de cerca de 700 estações automáticas do INMET, processados em um pipeline reprodutível em camadas. Uma auditoria dos dados identificou e corrigiu defeitos nos arquivos públicos, entre eles um deslocamento no mapeamento das colunas do INMET que fazia a variável tratada como umidade corresponder ao ponto de orvalho, e mostrou que parte dos resultados preliminares era artefato desses defeitos. Três variantes do XGBoost, com e sem variáveis climáticas, foram comparadas a previsores ingênuos em um teste de 2024 a 2026. O modelo com clima alcançou $R^2$ de 0,87 em escala logarítmica, acima de todos os previsores ingênuos, e superou a persistência em todas as métricas em 2025 e 2026; na epidemia de 2024, porém, antecipou o momento do pico, mas previu cerca de metade das notificações, e a persistência foi superior na escala original. O clima trouxe ganho pequeno, mas estatisticamente consistente, presente em todas as janelas de validação temporal e em cenários sem dados epidemiológicos recentes, e sozinho explicou 39\% da variância em escala logarítmica. Para a classificação do nível de risco, a persistência superou os classificadores dedicados. Os resultados foram disponibilizados em uma API REST e em painéis de visualização.

     \vspace{\onelineskip}

     \noindent
     \textbf{Palavras-chave}: Dengue. Previsão. Aprendizado de máquina. XGBoost. Clima. Vigilância epidemiológica.
\end{resumo}
